% Section 10.2, from lectures/L10/appendix/b_the_budget_and_the_loop.md.

\section{The budget, the missing operating point, and the loop}
\label{c10:app:b}

Where the gain goes, why the amplifier cannot be biased, and what closing the loop buys.

\subsection{Two stages of 66 dB give 120, not 131}
\label{c10:sec:b1}
\bookfigure{lectures/L10/appendix/images/gain_budget.png}{A waterfall chart of the gain budget in decibels: the differential pair at plus 65.7, small losses to the buffer, the common-emitter stage at plus 65.7, a loss of 10.6 to the output stage, and a total of 120 decibels.}{c10:fig:gainbudget}

\begin{narrowtable}{@{}ll@{}}
  \thead{} & \thead{dB} \\
  \midrule
  Stage 1, unloaded & +65.7 \\
  Loaded by the buffer & -0.1 \\
  The buffer itself & -0.3 \\
  Stage 3, unloaded & +65.7 \\
  Loaded by the output stage & \textbf{-10.6} \\
  The output follower itself & -0.5 \\
  \textbf{Open loop} & \textbf{119.9} \\
\end{narrowtable}

\textbf{11.5 dB of the 131.4 the two gain stages would give is lost to loading,} and not evenly: 10.6 is one number, stage 3 driving the output stage.

\textbf{This is Chapter~\ref{c1:ch:circuits}'s loading arithmetic for the last time,} the sixth appearance of one subtraction. \textbf{An amplifier's gain budget is not the product of its stage gains, and a budget computed without loading is wrong by a factor of 3.7 in the optimistic direction.}

\subsection{The two stages with no gain}
\label{c10:sec:b2}
Stage 2 has a voltage gain of 0.96 and stage 4 of 0.95. Together they are six of the eleven transistors and they amplify nothing. \textbf{Take each away and measure:}

\begin{booktable}{@{}Ll@{}}
  \thead{} & \thead{Open loop} \\
  \midrule
  As designed & \textbf{119.9 dB} \\
  With stage 1 driving stage 3 directly & 88.4 dB \\
  With a single follower instead of the output Darlington & 88.9 dB \\
\end{booktable}

\textbf{Each of the two stages with no gain is worth about 31 dB,} more than any two-transistor gain stage in the book could deliver. A stage that loses 0.3 dB of signal buys 31.5 dB of loop.

\textbf{That is the lesson the whole book has been building to.} Gain is easy: one transistor and a large load gives 66 dB. What is hard is \textbf{connecting two of them together}. Three transistors amplify, Q1, Q2 and Q5; six connect.

\subsection{The operating point that does not exist}
\label{c10:sec:b3}
Try to solve the amplifier for its DC operating point with the loop open. It will not work, and the reason is not the solver. The open-loop gain is 986,000 and the rails are $\pm 15$ V, so an input offset of

\[
  \frac{15}{986000} = 15\ \mu\text{V}
\]

puts the output hard against a rail. \textbf{A pair of transistors matched to 15 microvolts does not exist,} and neither does an amplifier whose open-loop output sits anywhere but a rail.

\textbf{So the DC solve either fails to converge or converges on a saturated output, and both are correct answers.} A solver reporting the output at $-15$ V has not failed.

\textbf{Which is why the stage gains have to be measured stage by stage.} Bias each stage alone, perturb its input, measure its output, multiply. The whole amplifier can only be solved with the loop closed, and then what you measure is the closed-loop gain of 20.

\textbf{And it is why stage 3 is undegenerated.} A designer looking at it in isolation would call it unbiasable and add an emitter resistor. Inside the loop it does not need one: the loop sets its operating point. \textbf{The stage is not a stage; it is part of a loop, and it cannot be judged alone.}

\subsection{Compensation, and the pole that is made dominant}
\label{c10:sec:b4}
Four stages give at least four poles, each up to 90 degrees of lag, and Chapter~\ref{c4:ch:feedback}~(Chapter~\ref{c4:ch:feedback}) established that 180 degrees with a loop gain above one is an oscillator. \textbf{The fix is not to make the amplifier fast, but to make it slow, deliberately, in one place.}

$C_c$ bridges stage 3's collector and base. It is a Miller capacitor (\secref{c7:sec:b5}), so the input sees $C_c(1 + 1923)$, which against that node's resistance puts a pole at a few tens of hertz. That resistance is not stage 3's 1.3 kilohm base resistance alone: the buffer's 72 ohm is in parallel, so the node sees 68 ohm. \textbf{The buffer is what makes $C_c$ large.} Every other pole is at hundreds of kilohertz or above.

\textbf{So by the time the other poles contribute phase, the loop gain is already below one.} The cost is bandwidth: the open-loop gain falls 20 dB per decade from that pole, and the closed-loop bandwidth is the gain-bandwidth product over the closed-loop gain. \textbf{This is Chapter~\ref{c7:ch:smallsignal}'s Miller effect used on purpose:} what ruined a common-emitter stage's bandwidth makes a four-stage amplifier usable.

\textbf{What this book does not do} is compute the other poles, so it cannot say what $C_c$ should be. That needs the capacitances of \secref{c7:sec:b8}.

\subsection{Closing the loop}
\label{c10:sec:b5}
Wrap a divider from the output back to the inverting input with a feedback fraction of $\beta_f = 1/20$, and Chapter~\ref{c4:ch:feedback}~(Chapter~\ref{c4:ch:feedback})'s arithmetic applies unchanged:

\begin{narrowtable}{@{}ll@{}}
  \thead{} & \thead{Value} \\
  \midrule
  Open-loop gain & 986,000 \\
  Feedback fraction & 1/20 \\
  Loop gain & 49,300 \\
  Closed-loop gain & 19.99959 \\
  Error from 20 & \textbf{0.002 per cent} \\
\end{narrowtable}

\textbf{And the thing that matters more than the error.} The output stage's input resistance goes as $h_{FE}^2$, which varies by a factor of a hundred across an ordinary beta spread:

\begin{narrowtable}{@{}lll@{}}
  \thead{$h_{FE}$} & \thead{Open loop} & \thead{Closed-loop gain} \\
  \midrule
  20 & 106.4 dB & 19.998096 \\
  50 & 119.9 dB & 19.999594 \\
  200 & 129.2 dB & 19.999862 \\
\end{narrowtable}

\textbf{The open-loop gain varies by a factor of 14. The closed-loop gain varies by nine parts in a hundred thousand.} That is the whole argument for feedback, on this book's own numbers. \secref{c8:sec:a5} said an output stage works inside a loop rather than on its own; this is the measurement.

\textbf{Output resistance} goes the same way. Open loop, three terms: the 0.22 ohm emitter resistor, the Darlington's $2r_e$ of 0.43 ohm, and stage 3's 50 kilohm over $h_{FE}^2$, which is 20 ohm and dominates. \textbf{20.6 ohm over one plus the loop gain: 0.42 milliohm.} Not a real number, since the wire to the loudspeaker limits the answer first, and saying so is the point of quoting it.

\subsection{What to build}
\label{c10:sec:b6}

\subsubsection{PNP support in \code{ael/net/netlist.hpp}}
The amplifier has four PNP devices and your netlist has none. Add a polarity:

\begin{cppcode}
enum class Polarity { Npn, Pnp };

void addBjt(Node collector, Node base, Node emitter, device::bjt::Parameters parameters = {},
            Polarity polarity = Polarity::Npn) noexcept;
\end{cppcode}

\textbf{A PNP is an NPN with every voltage and current negated:} the stamp you have, evaluated at $(-v_{BE}, -v_{BC})$ with the currents negated. Longer than fifteen lines and you are writing a second device model.

\subsubsection{\code{ael/report/amplifier.hpp}}
One structure, one function predicting every stage from Chapter~\ref{c7:ch:smallsignal} to Chapter~\ref{c9:ch:diffpair}'s closed forms, one formatter.

\begin{cppcode}
namespace ael::report
{
struct Design
{
    double tailCurrent{};       ///< Stage 1's tail, split between the two sides.
    double bufferCurrent{};     ///< Stage 2, the Darlington follower.
    double gainStageCurrent{};  ///< Stage 3, the common-emitter stage.
    double idleCurrent{};       ///< Stage 4's quiescent current, per half.
    double load{};              ///< The loudspeaker.
    double supply{};            ///< One rail, so the amplifier runs on plus and minus this.
    double beta{50.0};
    double earlyVoltage{100.0};
};

struct Stage
{
    std::string name{};
    double unloaded{};  ///< The stage on its own.
    double loadedBy{};  ///< The input resistance of the stage after it.
    double loaded{};    ///< What it delivers into that.
};

struct Budget
{
    std::vector<Stage> stages{};
};

/// The four stages, each loaded by the next stage's input resistance.
[[nodiscard]] Budget budget(const Design& design) noexcept;

/// The product of the loaded gains, and nothing else.
[[nodiscard]] double openLoopGain(const Design& design) noexcept;

/// L04's expression, applied to it.
[[nodiscard]] double closedLoopGain(const Design& design, double fraction) noexcept;

/// A table, one row per stage, predicted against measured.
[[nodiscard]] std::string format(const Budget& budget) noexcept;
}
\end{cppcode}

\textbf{The member names are part of the specification,} because the suite reads \code{budget.stages[0].loaded} and sets \code{design.tailCurrent} by name. \code{beta} and \code{earlyVoltage} are in \code{Design} rather than fixed, because \secref{c10:sec:b5}'s whole argument is what happens when they move.

\textbf{\code{budget} must call the functions from Chapter~\ref{c7:ch:smallsignal}, Chapter~\ref{c8:ch:follower} and Chapter~\ref{c9:ch:diffpair} rather than restating their formulas.} It is a composition, and if it contains an expression with $r_e$ in it then it has stopped being one. One shipped test checks that the stage-1 figure agrees with \code{ael::diffpair::mirrorGain} to machine precision.

\subsubsection{What good looks like}
About eighty lines, of which \code{format} is half. There is no physics in this component at all, which is what makes it the last one.

\subsection{What the capstone is for}
\label{c10:sec:b7}
The exercises end with one program that prints this:

\begin{console}
stage              predicted     measured    difference
1  input pair         1895.0       1837.2         -3.1 %
2  buffer                0.96         0.96        -0.1 %
3  voltage gain        570.4        548.6         -3.8 %
4  output                0.95         0.94        -0.9 %
open loop            986000       910000         -7.7 %
\end{console}

\textbf{The differences are the book,} and each has a name: the closed forms take $V_T$ as 26 mV where the device computes $kT/q$; they take $r_o$ as $V_A/I_C$ where the device gives $(V_A + V_{CE})/I_C$; beta is 50 in one and the transport model's incremental value in the other. \textbf{A reader who can account for every row has understood the ten chapters. A reader whose rows agree exactly has run the same code twice.}

\subsection{What this book is blind to}
\label{c10:sec:b8}
\begin{itemize}
  \item \textbf{Frequency response, beyond one pole.} No device capacitances, no phase margin, no Bode plot of a real amplifier. The largest omission, and where a second course would start.
  \item \textbf{Noise.} Not mentioned anywhere. It decides the input stage of every precision amplifier and most of the choices in \secref{c9:sec:b6}.
  \item \textbf{Distortion, quantitatively.} Where it comes from, never a number.
  \item \textbf{Temperature, beyond $V_{BE}$ drift.} No self-heating, no thermal models, no drift figures.
  \item \textbf{Integrated-circuit design.} Matching, layout, process variation and area are the subject next door; this book is discrete.
  \item \textbf{Anything above a few megahertz,} where transmission lines, package parasitics and electromagnetic compatibility take over from circuit theory.
\end{itemize}
\textbf{And one methodological blind spot.} Every number here comes from a model with eight parameters. Real devices have a hundred and real amplifiers are measured rather than computed. What the book teaches is how to \emph{predict}, which is worth having because it says what a measurement should be, and therefore when a measurement is telling you something.
